\documentclass[a4paper,12pt]{article}
\usepackage[utf8]{inputenc}
\usepackage[T1]{fontenc}
\usepackage[ngerman]{babel}
\usepackage{amsmath}
\usepackage{amssymb}
\usepackage{enumitem}
\usepackage{geometry}
\geometry{a4paper, margin=2.5cm}
\usepackage{parskip}

\title{Einsendeaufgaben -- Aufgabe 6.2\\[0.5em]
\large Modul 61112: Lineare Algebra}
\author{}
\date{}

\begin{document}
\maketitle

\section*{Aufgabe 1}

\[
A^2 =
\begin{pmatrix}
0 & 1 & 0 & 0 & 0 \\
1 & 1 & 0 & 0 & 1 \\
0 & 0 & 0 & 0 & 0 \\
5 & 5 & -1 & 0 & 5 \\
-1 & -2 & 0 & 0 & -1
\end{pmatrix}^2
=
\begin{pmatrix}
1 & 1 & 0 & 0 & 1 \\
0 & 0 & 0 & 0 & 0 \\
0 & 0 & 0 & 0 & 0 \\
0 & 0 & 0 & 0 & 0 \\
-1 & -1 & 0 & 0 & -1
\end{pmatrix}
\]

\[
A^3 =
\begin{pmatrix}
0 & 0 & 0 & 0 & 0 \\
0 & 0 & 0 & 0 & 0 \\
0 & 0 & 0 & 0 & 0 \\
0 & 0 & 0 & 0 & 0 \\
0 & 0 & 0 & 0 & 0
\end{pmatrix}
\]

Demnach ist $A$ nilpotent mit $\mathrm{Ni}(A) = 3$.

Nun bestimmen wir die Rangpartition:

\[
V_3 = \ker(A^3) = V, \quad \dim(V_3) = 5
\]
\[
V_2 = \ker(A^2) = \left\langle
\begin{pmatrix}1\\-1\\0\\0\\0\end{pmatrix},
\begin{pmatrix}0\\0\\1\\0\\0\end{pmatrix},
\begin{pmatrix}0\\0\\0\\1\\0\end{pmatrix},
\begin{pmatrix}1\\0\\0\\0\\-1\end{pmatrix}
\right\rangle, \quad \dim(V_2) = 4
\]
\[
V_1 = \ker(A) = \left\langle
\begin{pmatrix}0\\0\\0\\1\\0\end{pmatrix},
\begin{pmatrix}1\\0\\0\\0\\-1\end{pmatrix}
\right\rangle, \quad \dim(V_1) = 2, \text{ da}
\]

\[
\begin{pmatrix}
0 & 1 & 0 & 0 & 0 \\
1 & 1 & 0 & 0 & 1 \\
0 & 0 & 0 & 0 & 0 \\
5 & 5 & -1 & 0 & 5 \\
-1 & -2 & 0 & 0 & -1
\end{pmatrix}
\rightsquigarrow
\begin{pmatrix}
1 & 1 & 0 & 0 & 1 \\
0 & 1 & 0 & 0 & 0 \\
5 & 5 & -1 & 0 & 5 \\
-1 & -2 & 0 & 0 & -1 \\
0 & 0 & 0 & 0 & 0
\end{pmatrix}
\]
\[
\rightsquigarrow
\begin{pmatrix}
1 & 1 & 0 & 0 & 1 \\
0 & 1 & 0 & 0 & 0 \\
0 & 0 & -1 & 0 & 0 \\
0 & -1 & 0 & 0 & 0 \\
0 & 0 & 0 & 0 & 0
\end{pmatrix}
\rightsquigarrow
\begin{pmatrix}
1 & 1 & 0 & 0 & 1 \\
0 & 1 & 0 & 0 & 0 \\
0 & 0 & -1 & 0 & 0 \\
0 & 0 & 0 & 0 & 0 \\
0 & 0 & 0 & 0 & 0
\end{pmatrix}
\]

\[
p(A) = (2,\ 4-2,\ 5-4) = (2,2,1)
\]
\[
p(A)^* = (3,2)
\]

\[
\begin{array}{|c|c|c|}
\hline
v_{1,1} & v_{1,2} & v_{1,3} \\
\hline
v_{2,1} & v_{2,2} & \multicolumn{1}{c}{} \\
\cline{1-2}
\end{array}
\]

Wir wählen $v_{1,3} = \begin{pmatrix}1\\0\\0\\0\\0\end{pmatrix}$. Dann bestimmen wir

\[
v_{1,2} = A v_{1,3} = \begin{pmatrix}0\\1\\0\\5\\-1\end{pmatrix}
\quad\text{und}\quad
v_{1,1} = A^2 v_{1,3} = \begin{pmatrix}1\\0\\0\\0\\-1\end{pmatrix}.
\]

Wir wählen $v_{2,2} := \begin{pmatrix}0\\0\\1\\0\\0\end{pmatrix}$, da $\begin{pmatrix}0\\0\\1\\0\\0\end{pmatrix} \in V_2 \setminus V_1$

und $\langle v_{2,1}, v_{2,2}\rangle$ ein Komplement zu $V_1$ ist.

Wir bestimmen $v_{2,1} = \begin{pmatrix}0\\0\\0\\-1\\0\end{pmatrix}$.

Schließlich können wir $S$ aufstellen:

\[
S =
\begin{pmatrix}
1 & 0 & 1 & 0 & 0 \\
0 & 1 & 0 & 0 & 0 \\
0 & 0 & 0 & 0 & 1 \\
0 & 5 & 0 & -1 & 0 \\
-1 & -1 & 0 & 0 & 0
\end{pmatrix}
\]

Wir überprüfen, dass $S$ korrekt ist:

\[
\left(\begin{array}{ccccc|ccccc}
1 & 0 & 1 & 0 & 0 & 1 & 0 & 0 & 0 & 0 \\
0 & 1 & 0 & 0 & 0 & 0 & 1 & 0 & 0 & 0 \\
0 & 0 & 0 & 0 & 1 & 0 & 0 & 1 & 0 & 0 \\
0 & 5 & 0 & -1 & 0 & 0 & 0 & 0 & 1 & 0 \\
-1 & -1 & 0 & 0 & 0 & 0 & 0 & 0 & 0 & 1
\end{array}\right)
\]
\[
\rightsquigarrow
\left(\begin{array}{ccccc|ccccc}
1 & 0 & 1 & 0 & 0 & 1 & 0 & 0 & 0 & 0 \\
0 & 1 & 0 & 0 & 0 & 0 & 1 & 0 & 0 & 0 \\
0 & 0 & 0 & 0 & 1 & 0 & 0 & 1 & 0 & 0 \\
0 & 0 & 0 & -1 & 0 & 0 & -5 & 0 & 1 & 0 \\
0 & -1 & 1 & 0 & 0 & 1 & 0 & 0 & 0 & 1
\end{array}\right)
\]
\[
\rightsquigarrow
\left(\begin{array}{ccccc|ccccc}
1 & 0 & 1 & 0 & 0 & 1 & 0 & 0 & 0 & 0 \\
0 & 1 & 0 & 0 & 0 & 0 & 1 & 0 & 0 & 0 \\
0 & 0 & 0 & 0 & 1 & 0 & 0 & 1 & 0 & 0 \\
0 & 0 & 0 & 1 & 0 & 0 & 5 & 0 & -1 & 0 \\
0 & 0 & 1 & 0 & 0 & 1 & 1 & 0 & 0 & 1
\end{array}\right)
\]
\[
\rightsquigarrow
\left(\begin{array}{ccccc|ccccc}
1 & 0 & 0 & 0 & 0 & 0 & -1 & 0 & 0 & -1 \\
0 & 1 & 0 & 0 & 0 & 0 & 1 & 0 & 0 & 0 \\
0 & 0 & 0 & 0 & 1 & 0 & 0 & 1 & 0 & 0 \\
0 & 0 & 0 & 1 & 0 & 0 & 5 & 0 & -1 & 0 \\
0 & 0 & 1 & 0 & 0 & 1 & 1 & 0 & 0 & 1
\end{array}\right)
\]
\[
\rightsquigarrow
\left(\begin{array}{ccccc|ccccc}
1 & 0 & 0 & 0 & 0 & 0 & -1 & 0 & 0 & -1 \\
0 & 1 & 0 & 0 & 0 & 0 & 1 & 0 & 0 & 0 \\
0 & 0 & 1 & 0 & 0 & 1 & 1 & 0 & 0 & 1 \\
0 & 0 & 0 & 1 & 0 & 0 & 5 & 0 & -1 & 0 \\
0 & 0 & 0 & 0 & 1 & 0 & 0 & 1 & 0 & 0
\end{array}\right)
\]

\[
S S^{-1} =
\begin{pmatrix}
1 & 0 & 1 & 0 & 0 \\
0 & 1 & 0 & 0 & 0 \\
0 & 0 & 0 & 0 & 1 \\
0 & 5 & 0 & -1 & 0 \\
-1 & -1 & 0 & 0 & 0
\end{pmatrix}
\begin{pmatrix}
0 & -1 & 0 & 0 & -1 \\
0 & 1 & 0 & 0 & 0 \\
1 & 1 & 0 & 0 & 1 \\
0 & 5 & 0 & -1 & 0 \\
0 & 0 & 1 & 0 & 0
\end{pmatrix}
\]
\[
=
\begin{pmatrix}
1 & 0 & 0 & 0 & 0 \\
0 & 1 & 0 & 0 & 0 \\
0 & 0 & 1 & 0 & 0 \\
0 & 0 & 0 & 1 & 0 \\
0 & 0 & 0 & 0 & 1
\end{pmatrix}
\]

\[
S^{-1} A S = S^{-1}
\begin{pmatrix}
0 & 1 & 0 & 0 & 0 \\
1 & 1 & 0 & 0 & 1 \\
0 & 0 & 0 & 0 & 0 \\
5 & 5 & -1 & 0 & 5 \\
-1 & -2 & 0 & 0 & -1
\end{pmatrix}
\begin{pmatrix}
1 & 0 & 1 & 0 & 0 \\
0 & 1 & 0 & 0 & 0 \\
0 & 0 & 0 & 0 & 1 \\
0 & 5 & 0 & -1 & 0 \\
-1 & -1 & 0 & 0 & 0
\end{pmatrix}
\]
\[
=
\begin{pmatrix}
0 & -1 & 0 & 0 & -1 \\
0 & 1 & 0 & 0 & 0 \\
1 & 1 & 0 & 0 & 1 \\
0 & 5 & 0 & -1 & 0 \\
0 & 0 & 1 & 0 & 0
\end{pmatrix}
\begin{pmatrix}
0 & 1 & 0 & 0 & 0 \\
0 & 0 & 1 & 0 & 0 \\
0 & 0 & 0 & 0 & 0 \\
0 & 0 & 5 & 0 & -1 \\
0 & -1 & -1 & 0 & 0
\end{pmatrix}
\]
\[
=
\begin{pmatrix}
0 & 1 & 0 & 0 & 0 \\
0 & 0 & 1 & 0 & 0 \\
0 & 0 & 0 & 0 & 0 \\
0 & 0 & 0 & 0 & 1 \\
0 & 0 & 0 & 0 & 0
\end{pmatrix}
= \mathcal{N}\bigl(p(A)^*\bigr)
\]

\end{document}
